%%%%%%%%%%%%%%%%%%%%%%%%%%%%%%%%%%%%%%%%%%%%%%%%%%%%%%%%%%%%%%%%%%%%%%%%%%%%%%%
% Author:  Pablo Alvarado
%
% Área Académica de Ingeniería en Computadores
% Instituto Tecnológico de Costa Rica
%
% Tesis de Licenciatura
% 
% Phone:   +506 2550 2495
% email:   palvarado@tec.ac.cr
%
%%%%%%%%%%%%%%%%%%%%%%%%%%%%%%%%%%%%%%%%%%%%%%%%%%%%%%%%%%%%%%%%%%%%%%%%%%%%%%%

% \documentclass is book

% If you want a printable two-side version of the thesis
%\documentclass[12pt,twoside,letterpaper]{book}

% If you want an electronic-only version of the thesis, do it one-sided
\documentclass[12pt,oneside,letterpaper]{book}

\usepackage[utf8]{inputenc}
\usepackage{ifthen}                     % provide if-then-else operators
\usepackage[english]{babel}

% --------------------------------------------------------------------------
% Global variables required in document formatting
% --------------------------------------------------------------------------
%
% BOOK MODE
%
\newboolean{bookmode}                  % boolean used to control book format
% Ensure that only one of the next two lines is active:
\setboolean{bookmode}{true}           % turn book mode on
%\setboolean{bookmode}{false}           % turn book mode off

%
% DRAFT MODE
%   The draft mode activates the TODO index and some "draft" markings all
%   around.  Ensure you set it to false for the final version!!
%
\newboolean{draftmode}                  % boolean used to control draft-mode
% Ensure that only one of the next two lines is active:
%\setboolean{draftmode}{true}            % turn draft mode on
\setboolean{draftmode}{true}           % turn draft mode off
% --------------------------------------------------------------------------
%
% GENERAL AUTHOR, TITLE AND KEYWORDS
%
% Nombre del Estudiante
\newcommand{\thesisAuthorAddress}{el señor} %% el señor / la señora / ...
\newcommand{\thesisAuthor}{Luis Alonso Murillo Rojas} %% Nombre completo
\newcommand{\thesisAuthorShort}{L.~Murillo}  %% Nombre corto para pie de página
\newcommand{\thesisAuthorTECID}{201021497} %% Carné

% Título de la tesis
\newcommand{\thesisTitle}{Diseño de un aparato que produce energía
  infinita y elimina el cambio climático por medio del uso de
  circuitos analógicos controlados por una arquitectura RISC V}

% Keywords
\newcommand{\thesisKeywords}{Tomato, Deep Learning, Anomaly detection, Autoencoders, Generative Adversarial Networks}

% Descripción de la editorial
\newcommand{\boxeditorial}{%
}

% --------------------------------------------------------------------------

% include all packages and define all required general macros
%%%%%%%%%%%%%%%%%%%%%%%%%%%%%%%%%%%%%%%%%%%%%%%%%%%%%%%%%%%%%%%%%%%%%%%%%%%%%%%
% Author:  Pablo Alvarado
%
% Escuela de Electrónica
% Instituto Tecnológico de Costa Rica
%
% Phone:   +506 550 2106
% Fax:     +506 591 6629
% email:   palvarado@ietec.org
%
% $Id: macros.tex 1497 2010-08-09 17:04:26Z palvarado $
%
%%%%%%%%%%%%%%%%%%%%%%%%%%%%%%%%%%%%%%%%%%%%%%%%%%%%%%%%%%%%%%%%%%%%%%%%%%%%%%%

% Configuration of the exercises package, which is used to collect all
% problems and answers in the document.
\usepackage[exercisedelayed,answerdelayed,lastexercise]{exercise}
\renewcounter{Exercise}[chapter]
\renewcommand{\ExerciseName}{Problema}
\renewcommand{\theExercise}{\thechapter.\arabic{Exercise}}
\newcommand{\ExerciseLabel}{Exercise.\theExercise}
\renewcommand{\ExerciseHeader}%
{\textbf{\ExerciseName\ \theExercise.\ \ExerciseHeaderTitle\ }}
\renewcommand{\AnswerHeader}%
{\textbf{\ExerciseName\ \theExercise.\ }}


\usepackage{ifpdf}

% Command to change between draft or release mode:
\newcommand{\ifdraft}[2]{\ifthenelse{\boolean{draftmode}}{#1}{#2}}
% Command to change between draft or release mode:
\newcommand{\ifbook}[2]{\ifthenelse{\boolean{bookmode}}{#1}{#2}}

% include all required packages here
\usepackage{csquotes}                     % recommended for biblatex
\usepackage[spanish]{babel}               % supports english, but default is 
                                          % spanish...
% \newcommand*{\SelectSpanish}{%          % well, the last line indeed selects
%   \hyphenrules{spanish}%                % english over spanish, but with this
%   \languageshorthands{spanish}%         % command we turn it around.
%   \captionsspanish                      % The reason: hyperref has some
%   \datespanish                          % problems with the spanish babel,
% }                                       % so we use some trick here so that it
% \AtBeginDocument{\SelectSpanish}        % thinks it is english.

\usepackage{makeidx}                    % to create index file

\ifdraft{%
  %\usepackage[refpage]{nomencl}        % Use to easily administrate the list
 \usepackage{nomencl}                   % of symbols
}{%
 \usepackage{nomencl}
}
%\usepackage{times}                     % replace latex pk fonts with ps type I
                                        % don't forget to use dvips -D600 -Pcmz
                                        % to ensure Type I fonts!
\usepackage{amsmath}
\usepackage{amssymb,amstext}            % AMS-math and symbols package
\usepackage{mathrsfs}                   % Calygraphic fonts for transforms
\usepackage{array}                      % extensions to tabular environment
\usepackage{longtable}                  % supports extraordinary long tables
\usepackage{tabularx}                   % supports tables with fixed width
\usepackage{afterpage}                  % put something only after the page
\usepackage{multirow}                   % supports multiple row grouping in 
                                        % tables
\usepackage{multicol}                   % multiple columns environments
%\usepackage{paralist}                  % a few enumeration settings (old)
\usepackage{enumitem}                   % better enumeration with paralist 
                                        % equivalents as follows:

\newlist{compactitem}{itemize}{3}
\setlist[compactitem]{topsep=0pt,partopsep=0pt,itemsep=0pt,parsep=0pt}
\setlist[compactitem,1]{label=\textbullet}
\setlist[compactitem,2]{label=---}
\setlist[compactitem,3]{label=*}

\newlist{compactdesc}{description}{3}
\setlist[compactdesc]{topsep=0pt,partopsep=0pt,itemsep=0pt,parsep=0pt}

\newlist{compactenum}{enumerate}{3}
\setlist[compactenum]{topsep=0pt,partopsep=0pt,itemsep=0pt,parsep=0pt}

\usepackage{icomma}                     % decimal comma in math mode

\usepackage{bold-extra}

\usepackage[format=hang,%
            font=small,%
            labelfont=bf]{caption}      % nicer figure captions
%\usepackage{sty/ftcap}                 % switch \abovecaptionskip and
%                                       % \belowcaptionskip for tables, in 
%                                       % order to avoid the caption to be
%                                       % too near to the table itself
% locally added packages
\usepackage{float}                      % really place figures "here" (H)
\usepackage{booktabs}                   % book type tabulars

% the own style with options depending on the draft mode
\ifdraft{%
\usepackage[todo]{sty/tecStyle}         % some command definitions
                                        % options [todo] todo-index
}{%
\usepackage{sty/tecStyle}               % some command definitions
                                        % options [todo] todo-index
}


%% fix the title for examples
\renewcommand{\examplelistname}{Índice de ejemplos}
\renewcommand{\examplename}{Ejemplo}

%% define some command to cope with the tribunal names

%% Lector I

\newcommand{\nameLectorI}{$<$\emph{Use setLectorI in main.tex}$>$}
\newcommand{\genderLectorI}{$<$\emph{Use setLectorI in main.tex}$>$}

\newcommand{\setLectorI}[1][M]{%
  \ifthenelse{\equal{#1}{F}}{%
    \renewcommand{\genderLectorI}{Profesora Lectora}%
  }{%
    \renewcommand{\genderLectorI}{Profesor Lector}
  }
  \lectorIRelay
}

\newcommand{\lectorIRelay}[1]{%
  \renewcommand{\nameLectorI}{#1}
}

%% Lector II

\newcommand{\nameLectorII}{$<$\emph{Use setLectorII in main.tex}$>$}
\newcommand{\genderLectorII}{$<$\emph{Use setLectorII in main.tex}$>$}

\newcommand{\setLectorII}[1][M]{%
  \ifthenelse{\equal{#1}{F}}{%
    \renewcommand{\genderLectorII}{Profesora Lectora}%
  }{%
    \renewcommand{\genderLectorII}{Profesor Lector}
  }
  \lectorIIRelay
}

\newcommand{\lectorIIRelay}[1]{%
  \renewcommand{\nameLectorII}{#1}
}

%% Asesor

\newcommand{\nameAsesor}{$<$\emph{Use setAsesor in main.tex}$>$}
\newcommand{\genderAsesor}{$<$\emph{Use setAsesor in main.tex}$>$}

\newcommand{\setAsesor}[1][M]{%
  \ifthenelse{\equal{#1}{F}}{%
    \renewcommand{\genderAsesor}{Profesora Asesora}%
  }{%
    \renewcommand{\genderAsesor}{Profesor Asesor}
  }
  \asesorRelay
}

\newcommand{\asesorRelay}[1]{%
  \renewcommand{\nameAsesor}{#1}
}





\usepackage{url}                        % allows linebreaks at certain
                                        % characters or combinations of 
                                        % characters for URLs

\usepackage[nottoc]{tocbibind}          % Fix the hyperrefs to TOC,TOF, etc.
                                        % and ensure that they appear all in 
                                        % the Table of Contents

% For pdflatex
% - The hyperref package should always be loaded last, since it has to
%   overwrite some of the commands.
% - The package subfigure caused that the pagebackrefs and index refs were set
%   incorrectly.

\ifpdf
%
% final / draft document options
\usepackage{graphicx}                   % for inserting pdf-graphics.
                                        % options final / draft
\ifdraft{%
% Use biber/biblatex
\usepackage[backend=biber,
            style=ieee,
            sorting=nyt,
            backref=true,
           ]{biblatex}

\usepackage[%pdftex,%
            naturalnames=true,
            linktocpage,
            hyperindex,
            colorlinks,
            urlcolor=dkred,          %\href to external url
            filecolor=dkmagenta,     %\href to local file
            linkcolor=dkred,         %\ref and \pageref
            citecolor=dkgreen,       %\cite
            plainpages=false,
            pdfpagelabels,
            pdfpagemode=UseOutlines, % means use bookmarks (None,UseOutlines)
            % bookmarksopen=false,   % would show the whole hierarchy if true
            bookmarksnumbered=true,
            pdfpagelayout=OneColumn, % SinglePage,OneColumn,TwoColumnLeft,...
            pdfview=FitH, % FitB,FitBH,FitBV,Fit,FitH,FitV
            pdfstartview=FitH, % FitB,FitBH,FitBV,Fit,FitH,FitV
            ]{hyperref}
}{%
% Use biber/biblatex
\usepackage[backend=biber,
            style=ieee,
            sorting=nyt
           ]{biblatex}

\usepackage[%pdftex,%
            naturalnames=true,
            linktocpage,hyperindex,
            colorlinks,
            urlcolor=dkred,          %\href to external url
            filecolor=dkmagenta,     %\href to local file
            linkcolor=dkred,         %\ref and \pageref
            citecolor=dkgreen,       %\cite
            plainpages=false,
            pdfpagelabels,
            pdfpagemode=UseOutlines, % means use bookmarks (None,UseOutlines)
            % bookmarksopen=false,   % open the whole hierarchy if true!
            bookmarksnumbered=true,
            pdfpagelayout=OneColumn, % SinglePage,OneColumn,TwoColumnLeft,...
            pdfview=FitH, % FitB,FitBH,FitBV,Fit,FitH,FitV
            pdfstartview=FitH, % FitB,FitBH,FitBV,Fit,FitH,FitV
            ]{hyperref}
}


%
% Ensure that the links of the images point to the top of the images and not
% to the caption
%
\usepackage[figure]{hypcap}

% %
% % Ensure that pdfLaTeX do the same spacing as LaTeX
% %
\pdfadjustspacing=1 
% %
\else   % i.e. if not pdf

\usepackage[active]{srcltx}             % insert links into the dvi to jump
\usepackage{graphicx}                   % for inserting eps-graphics.
                                        % options final / draft
                                        % into the sources directly.
\ifdraft{%
% Use biber/biblatex
\usepackage[backend=biber,
            style=ieee,
            sorting=nyt,
            backref=true,
           ]{biblatex}

\usepackage[ps2pdf,%
            % plainpages=false,
            linktocpage,
            hyperindex,
            % pdfpagelabels,
            pdfpagemode=UseOutlines,
            pdfstartview=FitH]{hyperref}
}{%
% Use biber/biblatex
\usepackage[backend=biber,
            style=ieee,
            sorting=nyt
           ]{biblatex}
\usepackage[ps2pdf,%
            % plainpages=false,
            linktocpage,
            hyperindex,
            % pdfpagelabels,
            pdfpagemode=UseOutlines,
            pdfstartview=FitH]{hyperref}
}

%\usepackage[ps2pdf]{hyperref}

\fi  % end of if pdf or not

% --------------------------------------------------------------------------

% Allow the use of international characters
\AtBeginDocument{%
  \hypersetup{%
             pdftitle={\thesisTitle},%
             pdfsubject={Tesis de Licenciatura},%
             pdfauthor={\thesisAuthor},%
             pdfkeywords={\thesisKeywords}
            }%
}


%\usepackage{sty/algorithmic}            % algorithmic environment


\usepackage{rotating}                   % allow block rotation


%%%%%%%%%%%%%%%%%%%%%%%%%%%%%%%%%%%%%%%%%%%%%%%%%%%%%%%%%%%%%%%%%%%%%%%%%%%%%%%

%\sloppy

%
% Some own font definitions
%
\DeclareMathAlphabet{\mathpzc}{OT1}{pzc}{m}{it}
\DeclareMathAlphabet{\mathpss}{OT1}{cmss}{m}{sl}

%
% page layout
%

\usepackage{vmargin}
\setpapersize{USletter}

% For letter-paper printing
\setmarginsrb{33mm}{8mm}{23mm}{7mm}{15pt}{15pt}{7mm}{12mm}
%\setlength{\headheight}{15pt}         % fancy headers wanted this

%
% Fraction of Float Object / Text
%

\renewcommand{\topfraction}{0.95}       % how much of top of page should be 
                                        % allowed to be float object?
\renewcommand{\bottomfraction}{0.95}    % how much of bottom of page should be
                                        % allowed to be float object?
\renewcommand{\textfraction}{0.05}      % how much of page must be text?

\usepackage{fancyhdr}                   % fancy page headers

\usepackage{lastpage}

%
% header and footer layout (needs package fancyhdr)
%
\newcommand{\copyrightfooter}{\tiny{\copyright \the\year ---
    \thesisAuthorShort \qquad Uso exclusivo ITCR}}
%
\newcommand{\draftfoot}%
  {\ifdraft{\textcolor{dkblue}{\tiny\textsl{Borrador: \today}}{}}
           {}
}

\pagestyle{fancy}
\renewcommand{\chaptermark}[1]{\markboth{{\small
    \thechapter\hspace*{1mm}#1}}{}}
\renewcommand{\sectionmark}[1]{\markright{{\small
    \thesection\hspace*{1mm}#1}}{}}
%\lhead[{\small\textsc\Roman{\thepage}}]{\fancyplain{}%
\lhead[{\small\thepage}]{\fancyplain{}%
        {{\slshape \small\nouppercase{\leftmark}}}}
\chead[]{}
\rhead[\fancyplain{}%
%        {{\slshape \small\nouppercase{\rightmark}}}]{{\small\textsc\Roman{\thepage}}}
        {{\slshape \small\nouppercase{\rightmark}}}]{{\small\thepage}}
\lfoot[]{\draftfoot}
\ifbook{%
  \cfoot[]{}
}{
  \cfoot[\copyrightfooter]{\copyrightfooter}
}
\rfoot[\draftfoot]{}
\renewcommand{\headrulewidth}{0.5pt}
\renewcommand{\footrulewidth}{0pt}

%
% Caption style for tables
% Requires the packages caption2 and ftcap
% (caption2 required this, but is obsolete now:)
%
%\newcaptionstyle{tablecaptionstyle}{%
%  \renewcommand\captionlabelfont{\normalsize\bf}
%  \renewcommand\captionfont{\normalsize}
%  \usecaptionstyle{hang}%
%}

% For caption v3:
\captionsetup[table]{position=top,format=hang,textfont={normalsize},labelfont={normalsize,bf}}

\newcommand{\tablecaption}[2][foo]{%
  \ifthenelse{\equal{#1}{foo}}{%
    %\captionstyle{tablecaptionstyle}%
    \caption{#2}%
  }
  {%
    %\captionstyle{tablecaptionstyle}%
    \caption[#1]{#2}%
  }
}
\addto\extrasspanish{\renewcommand{\tablename}{Tabla}}
\addto\extrasspanish{\renewcommand{\listtablename}{\'Indice de tablas}}

%
% paragraph layout
%
\renewcommand{\baselinestretch}{1.1}    % line spacing
\parindent0em                           % indentation width of first line
\parskip1.3ex                           % space between paragraphs

%
% document consists of
% chapter - section - subsection - subsubsection - paragraph - subparagraph
%
\setcounter{secnumdepth}{2}             % depth of section numbering
\setcounter{tocdepth}{2}                % depth of table of contents

% For biblatex
\addbibresource{literatura.bib}

%
% prepares index from entries like \index{word} or \index{group!word}.
% don't forget to call "makeindex filename" for final index generation.
%
\makeindex                            %% for package makeidx.sty
%\newindex{default}{idx}{ind}{Index}  %% for package index.sty

\newcommand{\octave}{GNU/Octave}


%
% prepares notation or nomenclature 
%
%\makeglossary
\makenomenclature

%%% Local Variables: 
%%% mode: latex
%%% TeX-master: "main"
%%% End: 


% Tribunal Evaluador
%
% Indique los nombres de los lectores y asesor
% El parámetro opcional es [M]asculino o [F]emenino
\setLectorI[F]{Dra.\, María del Pilar Pérez Fernández}
\setLectorII{M.\,Sc.\,Juan Pérez Hernández}
\setAsesor{Ing.\,Albert Einstein Sánchez}

% allow equations to be splitted (breaked) into several pages
\allowdisplaybreaks[3]

% --------------------------------------------------------------------------
\begin{document}
  % where to look for graphics
  \graphicspath{{./}{./fig/}}

  \pagenumbering{alph}
  % fix some terms not activated due to the bug of hyperref with spanish.
  \renewcommand{\tablename}{Table}
  \renewcommand{\listtablename}{\'Table index}
  \renewcommand{\examplesolution}{Solution}
  \pagestyle{empty}

  % select one of the following titlepages
  %\include{titlepage_licel_es} % Titlepage in Spanish

  %%% ---------------------------------------------------------------------------
%% titlepage_licce_es.tex
%%
%% Title page
%%
%% $Id: titlepage.tex 1452 2010-07-07 00:55:16Z palvarado $
%% ---------------------------------------------------------------------------

\thispagestyle{empty} 

\begin{center}

Tecnológico de Costa Rica

\par\vspace{1ex}

Área Académica de Ingeniería en Computadores

\par\vspace{1ex}

Programa de Licenciatura de Ingeniería en Computadores

\par\vspace{20mm}

\includegraphics[width=60mm]{Firma_TEC-4}

\par\vspace*{\fill}

{\large\bf{\thesisTitle}}

\par\vspace*{\fill}

Informe de Trabajo Final de Graduación para optar por el título de

Ingeniero en Computadores con el grado académico de Licenciatura

\par\vspace{20mm}

\thesisAuthor

\vspace*{\fill}

\ifdraft{%
{Borrador de \today}
}{
Cartago, 7 de marzo, 2018
}
\end{center}
\newpage 
\cleardoublepage 


%%% Local Variables: 
%%% mode: latex
%%% TeX-master: "main"
%%% End: 
  % Titlepage in Spanish
  %%% ---------------------------------------------------------------------------
%% titlepage_licce_es.tex
%%
%% Title page
%%
%% $Id: titlepage.tex 1452 2010-07-07 00:55:16Z palvarado $
%% ---------------------------------------------------------------------------

\thispagestyle{empty} 

\begin{center}

Tecnológico de Costa Rica

\par\vspace{1ex}

Computer Engineering Department

\par\vspace{1ex}

Licentiate Degree Program in Computer Engineering

\par\vspace{20mm}

\includegraphics[width=60mm]{Firma_TEC-4}

\par\vspace*{\fill}

{\large\bf{\thesisTitle}}

\par\vspace*{\fill}

Final report submitted in partial fulfillment of the requirements
for the degree of Licentiate in Computer Engineering

\par\vspace{20mm}

\thesisAuthor

\vspace*{\fill}

\ifdraft{%
{Draft \today}
}{
Cartago, March 7th, 2018
}
\end{center}
\newpage 
\cleardoublepage 


%%% Local Variables: 
%%% mode: latex
%%% TeX-master: "main"
%%% End: 
 % Titlepage in English (only if thesis is in En)
  
  %%% ---------------------------------------------------------------------------
%% titlepage_msc_es.tex
%%
%% Title page
%%
%% $Id: titlepage.tex 1452 2010-07-07 00:55:16Z palvarado $
%% ---------------------------------------------------------------------------

\thispagestyle{empty} 

\begin{center}

Tecnológico de Costa Rica

\par\vspace{1ex}

Escuela de Ingeniería Electrónica

\par\vspace{20mm}

\includegraphics[height=60mm]{Firma_TEC-4}

\par\vspace*{\fill}

{\large\bf{\thesisTitle}}

\par\vspace*{\fill}

Documento de tesis sometido a consideración para optar por el grado
académico de Maestría en Electrónica con Énfasis en
%
%Sistemas Embebidos
Procesamiento Digital de Señales
%Microelectrónica
%Sistemas Microelectromecánicos

\par\vspace{20mm}

\thesisAuthor

\vspace*{\fill}

\ifdraft{%
{Borrador de \today}
}{
Cartago, 19 de noviembre, 2013
}
\end{center}
\newpage 
\cleardoublepage 


%%% Local Variables: 
%%% mode: latex
%%% TeX-master: "main"
%%% End: 
   % Titlepage in Spanish
  %% ---------------------------------------------------------------------------
%% titlepage.tex
%%
%% Title page
%%
%% $Id: titlepage.tex 1452 2010-07-07 00:55:16Z palvarado $
%% ---------------------------------------------------------------------------

\thispagestyle{empty} 

\begin{center}

Tecnológico de Costa Rica

\par\vspace{1ex}

Escuela de Ingeniería Electrónica

\par\vspace{20mm}

\includegraphics[height=60mm]{Firma_TEC-4}

\par\vspace*{\fill}

{\large\bf{\thesisTitle}}

\par\vspace*{\fill}

A thesis submitted in partial fulfillment of the requirements for the
degree of
%
Master of Science in Electronics, Major in 
%
%Embedded Systems
Digital Signal Processing
%Microelectronics
%Microelectromechanical systems

\par\vspace{20mm}

\thesisAuthor

\vspace*{\fill}

\ifdraft{%
{Draft \today}
}{
Cartago, November 19th, 2017
}
\end{center}
\newpage 
\cleardoublepage 


%%% Local Variables: 
%%% mode: latex
%%% TeX-master: "main"
%%% End: 
   % Titlepage in English (only if thesis is in En)

  \thispagestyle{empty}

\rule{10mm}{0pt}

\vfill

Declaro que el presente documento de tesis ha sido realizado enteramente
por mi persona, utilizando y aplicando literatura referente al tema e
introduciendo conocimientos y resultados experimentales propios.

En los casos en que he utilizado bibliografía he procedido a indicar las
fuentes mediante las respectivas citas bibliográficas.  En consecuencia,
asumo la responsabilidad total por el trabajo de tesis realizado y por
el contenido del presente documento.



\vspace*{8mm}

\begin{flushright}
  \thesisAuthor\par
  Cartago, \today\par
  Céd: 1-0123-0456
\end{flushright}

\cleardoublepage

%%% Local Variables: 
%%% mode: latex
%%% TeX-master: "main"
%%% End: 

  
  %% -------------------------------------------------
  %% Acta y hoja del tribunal
  
  %% Para la licenciatura en electrónica:
  %%% ESTE ARCHIVO DEBE ELIMINARSE DE LA VERSIÓN FINAL

\thispagestyle{empty}

\begin{center}
  \begin{tabular}{c}
    Instituto Tecnológico de Costa Rica \\
    Escuela de Ingeniería Electrónica \\
    Proyecto de Graduación \\
    Acta de Aprobación
  \end{tabular}
\end{center}

\vfill

\begin{center}
  \begin{tabular}{c}
    Defensa de Proyecto de Graduación \\
    Requisito para optar por el título de Ingeniero en Electrónica \\
    Grado Académico de Licenciatura
  \end{tabular}
\end{center}

\vfill

%% \thesisAuthorAddress, \thesisAuthor y \thesisTitle están en main.tex
El Tribunal Evaluador aprueba la defensa del proyecto de graduación
denominado \textsl{\thesisTitle}, realizado por
%
\thesisAuthorAddress\ \thesisAuthor\ %
%
y, hace constar que cumple con las normas
establecidas por la Escuela de Ingeniería Electrónica del Instituto
Tecnológico de Costa Rica.

\vfill

\begin{center}
 Miembros del Tribunal Evaluador
\end{center}

\vfill

\begin{center}
  \begin{tabularx}{\textwidth}{cXc}
    \rule{0.45\textwidth}{0.5pt} && \rule{0.45\textwidth}{0.5pt} \\
    \nameLectorI                 && \nameLectorII \\
    \genderLectorI               && \genderLectorII
  \end{tabularx}
  
  \vspace{10mm}

  \begin{tabular}{c}
    \rule{0.45\textwidth}{0.5pt} \\
    \nameAsesor \\
    \genderAsesor
  \end{tabular}
\end{center}

\vfill

\begin{center}
  Cartago, \today\par
\end{center}

\cleardoublepage

%%% Local Variables: 
%%% mode: latex
%%% TeX-master: "main"
%%% End: 
  % Remover en versión final
  %%% ESTE ARCHIVO DEBE ELIMINARSE DE LA VERSIÓN FINAL


\thispagestyle{empty}

\begin{center}
  \begin{tabular}{c}
    Instituto Tecnológico de Costa Rica \\
    Escuela de Ingeniería Electrónica \\
    Proyecto de Graduación \\
    Tribunal Evaluador \\
    Acta de Evaluación
  \end{tabular}
\end{center}

\vfill

\begin{center}
  \begin{tabular}{c}
    Defensa de Proyecto de Graduación \\
    Requisito para optar por el título de Ingeniero en Electrónica \\
    Grado Académico de Licenciatura
  \end{tabular}
\end{center}

\vfill

%% \thesisAuthorAddress, \thesisAuthor y \thesisTitle están en main.tex
\begin{center}

  Estudiante:%
  \qquad \textbf{\thesisAuthor}%
  \qquad Carné: \thesisAuthorTECID

  \vspace*{2ex}

  \setlength\tabcolsep{0pt}
  \begin{tabular}{p{.25\textwidth}p{.73\textwidth}}
    Nombre del proyecto: & \textsl{\thesisTitle}
  \end{tabular}
\end{center}
\vspace{5mm}

\vfill

Los miembros de este Tribunal hacen constar que este proyecto de
graduación ha sido aprobado y cumple con las normas establecidas por
la Escuela de Ingeniería Electrónica del Instituto Tecnológico de
Costa Rica y es merecedor de la siguiente calificación:

\vfill

\begin{center}
  Nota de Proyecto de Graduación: \rule{25mm}{0.5pt}
\end{center}

\vfill

\begin{center}
 Miembros del Tribunal Evaluador
\end{center}

\vfill

% Defina con \setLector* en main.tex los lectores y asesor
\begin{center}
  \begin{tabularx}{\textwidth}{cXc}
    \rule{0.45\textwidth}{0.5pt} && \rule{0.45\textwidth}{0.5pt} \\
    \nameLectorI                 && \nameLectorII \\
    \genderLectorI               && \genderLectorII
  \end{tabularx}
  
  \vspace{10mm}

  \begin{tabular}{c}
    \rule{0.45\textwidth}{0.5pt} \\
    \nameAsesor \\
    \genderAsesor
  \end{tabular}
\end{center}

\vfill

\begin{center}
  Cartago, \today\par
\end{center}

\cleardoublepage

%%% Local Variables: 
%%% mode: latex
%%% TeX-master: "main"
%%% End: 
      % Remover en versión final
  
  %% Para la maestría en electrónica:
  %% ESTE ARCHIVO DEBE ELIMINARSE DE LA VERSIÓN FINAL

\thispagestyle{empty}

\begin{center}
  \begin{tabular}{c}
    Instituto Tecnológico de Costa Rica \\
    Escuela de Ingeniería Electrónica \\
    Proyecto de Graduación \\
    Tesis de Maestría \\
    Tribunal Evaluador
  \end{tabular}
\end{center}

\vfill

Tesis de maestría defendida ante el presente Tribunal Evaluador como
requisito para optar por el grado académico de maestría, del Instituto
Tecnológico de Costa Rica.

\vfill

\vspace*{20mm}
\begin{center}
 Miembros del Tribunal
\end{center}
\vspace*{8mm}

\vfill

\begin{center}
  \begin{tabularx}{\textwidth}{cXc}
    \rule{0.45\textwidth}{0.5pt} && \rule{0.45\textwidth}{0.5pt} \\
    \nameLectorI                 && \nameLectorII \\
    \genderLectorI               && \genderLectorII
  \end{tabularx}
  
  \vspace{10mm}

  \begin{tabular}{c}
    \rule{0.45\textwidth}{0.5pt} \\
    \nameAsesor \\
    \genderAsesor
  \end{tabular}
\end{center}

\vfill


Los miembros de este Tribunal dan fe de que la presente tesis de maestría 
ha sido aprobada y cumple con las normas establecidas por la Escuela de
Ingeniería Electrónica.

\vfill

\begin{center}
  Cartago, \today\par
\end{center}

\cleardoublepage

%%% Local Variables: 
%%% mode: latex
%%% TeX-master: "main"
%%% End: 
  % Remover en versión final
  %% ESTE ARCHIVO DEBE ELIMINARSE DE LA VERSIÓN FINAL


\thispagestyle{empty}

\begin{center}
  \begin{tabular}{c}
    Instituto Tecnológico de Costa Rica \\
    Escuela de Ingeniería Electrónica \\
    Tesis de Maestría \\
    Acta de Evaluación
  \end{tabular}
\end{center}

\vfill

Tesis de maestría defendida ante el presente Tribunal Evaluador como
requisito para optar por el grado académico de maestría, del Instituto
Tecnológico de Costa Rica.

\vspace*{15mm}

%% Definir \thesisAuthor, \thesisTitle, etc. en archivo main.tex
\begin{center}
  Estudiante: \thesisAuthor
\end{center}

\vfill

\begin{center}
  Nombre del Proyecto: \thesisTitle}
\end{center}

\vspace*{20mm}
\begin{center}
 Miembros del Tribunal Evaluador
\end{center}
\vspace*{8mm}

\vfill

% Los nombres de lectores y asesor se definen en el archivo main.tex
\begin{center}
  \begin{tabularx}{\textwidth}{cXc}
    \rule{0.45\textwidth}{0.5pt} && \rule{0.45\textwidth}{0.5pt} \\
    \nameLectorI                 && \nameLectorII \\
    \genderLectorI               && \genderLectorII
  \end{tabularx}
  
  \vspace{10mm}

  \begin{tabular}{c}
    \rule{0.45\textwidth}{0.5pt} \\
    \nameAsesor \\
    \genderAsesor
  \end{tabular}
\end{center}

\vfill

Los miembros de este Tribunal dan fe de que la presente tesis de
maestría ha sido aprobada y cumple con las normas establecidas por la
Escuela de Ingeniería Electrónica.

\vfill

\begin{center}
  Nota final de la Tesis de Maestría: \rule{3cm}{0.5pt}
\end{center}
\vfill

\begin{center}
  Cartago, \today\par
\end{center}

\cleardoublepage

%%% Local Variables: 
%%% mode: latex
%%% TeX-master: "main"
%%% End: 
      % Remover en versión final
  %% -------------------------------------------------
  
  \chapter*{Resumen}
\thispagestyle{empty}

El resumen es la síntesis de lo que aparecerá en el tesis. Tiene que ser lo
suficientemente consiso y claro para que alguien que lo lea sepa qué esperar
del resto de la tesis si la leyera completamente. Puede concluir con palabras
clave, que son los temas principales tratados en el documento. El resumen queda
fuera de la numeración del resto de secciones.

No se acostumbra utilizar referencias bibliográficas, tablas, o figuras
en el resumen.

\bigskip

\textbf{Palabras clave:} \thesisKeywords

\clearpage
\chapter*{Abstract}
\thispagestyle{empty}

The same as before, but in English.

\bigskip

\textbf{Keywords:} word 1, word 2, 

\cleardoublepage

%%% Local Variables: 
%%% mode: latex
%%% TeX-master: "main"
%%% End: 

  \vspace*{0.4\textheight}
{\hfill{\Large{\emph{a Mariángel}}}}

  \chapter*{Agradecimientos}
\thispagestyle{empty}

El resultado de este trabajo no hubiese sido posible sin el apoyo de Thevenin,
Norton, Einstein y mi querido amigo Ohm.

\vspace*{1cm}

\thesisAuthor

Cartago, \today

\cleardoublepage

%%% Local Variables: 
%%% mode: latex
%%% TeX-master: "paMain"
%%% End: 


  %----------------------------------------------------------------------------
  \frontmatter
  %----------------------------------------------------------------------------
  \pagestyle{fancy}
  \pagenumbering{roman}

  \pdfbookmark[1]{General index}{General index}

  \parskip0ex                           % space between paragraphs

  \tableofcontents                      % Table of contents
  \listoffigures                        % List of figures
  \listoftables                         % List of tables

\ifdraft{%
  % todo's                              % TODOs
  \listoftodo
}{%
}

  %% ---------------------------------------------------------------------------
%% paNotation.tex
%%
%% Notation
%%
%% $Id: notation.tex 1467 2010-07-24 16:47:17Z palvarado $
%% ---------------------------------------------------------------------------

\newcommand{\nms}{\negmedspace}

%%
% Commands required for the nomenclature groups
%
% There are following prefix forms:
%  a   abbreviation    \syma[key]{symbol}{description}
%  g   general         \symg[key]{symbol}{description}
%%

\newcommand{\nmstyle}[1]{\large\item[\textbf{#1}]\normalsize}

\renewcommand{\nomgroup}[1]{%
  \ifthenelse{\equal{#1}{A}}{\nmstyle{Abreviaciones}\normalsize}{%
  \ifthenelse{\equal{#1}{G}}{\bigskip\nmstyle{Notación general}\normalsize}%
  }
}

\newcommand{\syma}[3][foo]{%
  \ifthenelse{\equal{#1}{foo}}%
  {\nomenclature[A#2\ ]{#2}{#3}}{\nomenclature[A#1\ ]{#2}{#3}}}
\newcommand{\symg}[3][foo]{%
  \ifthenelse{\equal{#1}{foo}}%
  {\nomenclature[G#2\ ]{#2}{#3}}{\nomenclature[G#1\ ]{#2}{#3}}}

%%
% Command definitions for localized symbol format definition
%%
\renewcommand{\Re}{\operatorname{Re}}
\renewcommand{\Im}{\operatorname{Im}}

\newcommand{\prt}[1]{\ensuremath{\mathcal{#1}}}         %% partitioning
\newcommand{\img}[1]{\ensuremath{\mathcal{#1}}}         %% image as a set
\newcommand{\reg}[1][R]{\ensuremath{\mathcal{#1}}}      %% region
\newcommand{\pred}[1]{\ensuremath{\mathrm{#1}}}         %% predicate
\newcommand{\operat}[2]{\mathcal{#1}\left\{#2\right\}}
\newcommand{\transf}[1]{\mathscr{#1}}
\newcommand{\fourier}[1]{\transf{F}\left\{#1\right\}}
\newcommand{\ifourier}[1]{\transf{F}^{-1}\left\{#1\right\}}
\newcommand{\laplace}[1]{\transf{L}\left\{#1\right\}}
\newcommand{\ulaplace}[1]{\transf{L}_u\left\{#1\right\}}
\newcommand{\blaplace}[1]{\transf{L}_b\left\{#1\right\}}
\newcommand{\ilaplace}[1]{\transf{L}^{-1}\left\{#1\right\}}
\newcommand{\ztrans}[1]{\transf{Z}\left\{#1\right\}}
\newcommand{\iztrans}[1]{\transf{Z}^{-1}\left\{#1\right\}}
\newcommand{\zutrans}[1]{\transf{Z}_u\left\{#1\right\}}
\newcommand{\exceq}{\ensuremath{\overset{!}{=}}}

\newcommand{\signum}{\operatorname{signum}}
\newcommand{\vct}[1]{\ensuremath{\underline{\mathbf{#1}}}}
\newcommand{\mat}[1]{\ensuremath{\mathbf{#1}}}
\newcommand{\vctmu}{\vct{\boldsymbol{\mu}}}
\newcommand{\vctzeta}{\vct{\boldsymbol{\zeta}}}
\newcommand{\vctpi}{\vct{\boldsymbol{\pi}}}
\newcommand{\vctvarphi}{\vct{\boldsymbol{\varphi}}}
\newcommand{\raum}[1]{\ensuremath{\mathbb{#1}}}
\newcommand{\matSigma}{\mat{\boldsymbol{\Sigma}}}
\newcommand{\matLambda}{\mat{\boldsymbol{\Lambda}}}
\newcommand{\matPsi}{\mat{\boldsymbol{\Psi}}}
\newcommand{\matPhi}{\mat{\boldsymbol{\Phi}}}
\newcommand{\row}[2]{\ensuremath{\mathbf{\underline{#1}_{#2(\cdot)}}}}
\newcommand{\col}[2]{\ensuremath{\mathbf{\underline{#1}_{(\cdot) #2}}}}
\newcommand{\seq}[1]{\ensuremath{#1}}
\newcommand{\set}[1]{\ensuremath{\mathcal{#1}}}
\newcommand{\gset}[1]{\ensuremath{#1}} %% set for greek symbols
\newcommand{\front}[1]{\widehat{\set{#1}}}
\newcommand{\setlambda}{\set{\boldsymbol{\lambda}}}
\newcommand{\klass}[1]{\ensuremath{\mathpss{#1}}}
\newcommand{\graph}[1]{\ensuremath{\mathsf{#1}}}
\newcommand{\lab}[1]{\ensuremath{\mathpss{L}(#1)}}
\newcommand{\myfrac}[2]{{\footnotesize #1/#2}}
\newcommand{\ifthenspc}{\rule{3mm}{0mm}}
\newcommand{\point}[1]{\ensuremath{\mathsf{#1}}}
\newcommand{\estim}[1]{\ensuremath{\hat{#1}}}
\newcommand{\numset}[1]{\ensuremath{\mathbb{#1}}}
\newcommand{\tuple}[1]{\ensuremath{\left\langle#1\right\rangle}}
\newcommand{\norm}[1]{\ensuremath{\left\lVert#1\right\rVert}}
\newcommand{\conj}[1]{\ensuremath{{{#1}^{\ast}}}}
\newcommand{\base}[1]{\set{#1}}
\newcommand{\zeron}[1]{\ensuremath{\underset{\uparrow}{#1}}}
\newcommand{\sysT}{\ensuremath{\mathcal{T}}}
\newcommand{\sys}[1]{\ensuremath{\sysT\left[#1\right]}}
%\newcommand{\sen}{\operatorname{sen}} % sinus in spanish (seno)
%\newcommand{\senh}{\operatorname{senh}} % sinus hiperbolicus in spanish (seno)
%\newcommand{\arcsen}{\operatorname{arcsen}} % arcus sinus hiperbolicus in spanish (arcoseno)
\newcommand{\sgn}{\operatorname{sgn}} % signus
\newcommand{\roc}{\text{ROC: }}

\newcommand{\code}[1]{\texttt{#1}}
\newcommand{\conv}{\ensuremath{\ast}}
\newcommand{\cconv}{\ensuremath{\;\,\text{\footnotesize{N}}\!\!\!\!\!\!\bigcirc}}
\newcommand{\Ln}{\operatorname{Ln}}
\newcommand{\rand}{\operatorname{rand}}
\newcommand{\sa}{\operatorname{sa}}
\newcommand{\senc}{\operatorname{senc}}
\newcommand{\si}{\operatorname{si}}


%% Natural, Integer and Real Numbers
\newcommand{\setA}{\ensuremath{\mathbb{A}}}
\newcommand{\setB}{\ensuremath{\mathrm{I\negthinspace B}}}
\newcommand{\setC}{\ensuremath{\mathbb{C}}}
\newcommand{\setD}{\ensuremath{\mathrm{I\negthinspace D}}}
\newcommand{\setE}{\ensuremath{\mathrm{I\negthinspace E}}}
\newcommand{\setF}{\ensuremath{\mathrm{I\negthinspace F}}}
\newcommand{\setG}{\ensuremath{\mathbb{G}}}
\newcommand{\setH}{\ensuremath{\mathrm{I\negthinspace H}}}
\newcommand{\setI}{\ensuremath{\mathbb{I}}}
\newcommand{\setJ}{\ensuremath{\mathbb{J}}}
\newcommand{\setK}{\ensuremath{\mathrm{I\negthinspace K}}}
\newcommand{\setL}{\ensuremath{\mathrm{I\negthinspace L}}}
\newcommand{\setM}{\ensuremath{\mathrm{I\negthinspace M}}}
\newcommand{\setN}{\ensuremath{\mathrm{I\negthinspace N}}}
\newcommand{\setO}{\ensuremath{\mathbb{O}}}
\newcommand{\setP}{\ensuremath{\mathrm{I\negthinspace P}}}
\newcommand{\setQ}{\ensuremath{\mathbb{Q}}}
\newcommand{\setR}{\ensuremath{\mathrm{I\negthinspace R}}}
\newcommand{\setS}{\ensuremath{\mathbb{S}}}
\newcommand{\setT}{\ensuremath{\mathbb{T}}}
\newcommand{\setU}{\ensuremath{\mathbb{U}}}
\newcommand{\setV}{\ensuremath{\mathbb{V}}}
\newcommand{\setW}{\ensuremath{\mathbb{W}}}
\newcommand{\setX}{\ensuremath{\mathbb{X}}}
\newcommand{\setY}{\ensuremath{\mathbb{Y}}}
\newcommand{\setZ}{\ensuremath{\mathbb{Z}}}


%%
% Multimap symbols
%
\newcommand{\ttoF}{\,\circ\!\negthickspace\longrightarrow\negthickspace\!\negthickspace\bullet\,}
\newcommand{\Ftot}{\,\bullet\negthickspace\!\negthickspace\longleftarrow\!\negthickspace\circ\,}
\newcommand{\ttoZ}{\ttoF}
\newcommand{\Ztot}{\Ftot}
\newcommand{\ttoZu}{\overset{z_u}{\ttoF}}
\newcommand{\Zutot}{\overset{z_u}{\Ftot}}
\newcommand{\vttoF}{\text{\begin{sideways}$\Ftot$\end{sideways}}}
\newcommand{\vFtot}{\text{\begin{sideways}$\ttoF$\end{sideways}}}
\newcommand{\vttoZ}{\vttoF}
\newcommand{\vZtot}{\vFtot}
\newcommand{\ttoDF}{\underset{N}{\ttoF}}
\newcommand{\DFtot}{\underset{N}{\Ftot}}

\newcommand{\thisis}[2]{\underset{#1}{\underbrace{#2}}}

%%% Local Variables:
%%% mode: latex
%%% TeX-master: "paMain"
%%% End:
                    % Notation
  %% ---------------------------------------------------------------------------
%% paNotation.tex
%%
%% Notation
%%
%% $Id: paNotation.tex,v 1.15 2004/03/30 05:55:59 alvarado Exp $
%% ---------------------------------------------------------------------------

\cleardoublepage
\renewcommand{\nomname}{List of symbols and abbreviations}
\markboth{\nomname}{\nomname}
\renewcommand{\nompreamble}{\addcontentsline{toc}{chapter}{\nomname}%
\setlength{\nomitemsep}{-\parsep}
\setlength{\itemsep}{10ex}
}

%%
% Símbolos en la notación general
% (es posible poner la declaración en el texto
%%

\symg[t]{$\sys{\cdot}$}{Transformation performed by a system}
\symg[yscalar]{$y$}{Scalar.}
\symg[zconjugado]{$\conj{z}$}{Conjugate Complex of $z$}
\symg[rcomplexreal]{$\Re(z)$ o $z_{\Re}$}{Real part of the complex number $z$}
\symg[icompleximag]{$\Im(z)$ o $z_{\Im}$}{Imaginary part of the complex number $z$}
\symg[jimaginario]{$j$}{$j=\sqrt{-1}$}
\symg[xvector]{$\vct{x}$}{Vector. \newline\hspace{1mm}%
  $\vct{x}=\left[ x_1 \; x_2 \; \ldots \; x_n \right]^T =
  \begin{bmatrix}
    x_1 \\ x_2 \\ \vdots \\ x_n
  \end{bmatrix}$}

\symg[amatrix]{$\mat{A}$}{Matrix. \newline\hspace{1mm}%
  $\mat{A} =
  \begin{bmatrix}
    a_{11} & a_{12} & \cdots & a_{1m}\\
    a_{21} & a_{22} & \cdots & a_{2m}\\
    \vdots & \vdots & \ddots & \vdots\\
    a_{n1} & a_{n2} & \cdots & a_{nm}\\
  \end{bmatrix}$}

\symg[C]{$\setC$}{Complex numbers set.}

%%
% Algunas abreviaciones
%%

\syma{MLP}{Multilayer perceptrons}
\syma{CNN}{Convolutional Neural Networks}
\syma{VAE}{Variational Autoencoder}
\syma{sVAE}{Spatial Variational Autoencoder}
\syma{GM-VAE}{Gaussian-Mixture Variational Autoencoder}
\syma{GAN}{Generative Adversarial Networks}
\syma{AAD}{Adversarial Anomaly Detector}
\syma{KL}{Kullback-Leibler}
\syma{MVN}{Matrix-variable normal distribution}
\syma{STD}{Standard deviation}

\printnomenclature[20mm]

%%% Local Variables:
%%% mode: latex
%%% TeX-master: "paMain"
%%% End:
                    % Abbreviation

  \parskip1.3ex                         % space between paragraphs

  %----------------------------------------------------------------------------
  \mainmatter
  %----------------------------------------------------------------------------
  % where to look for graphics
  \graphicspath{{./}{./fig/}}
  %\pagenumbering{arab}

  % Main files
  %% ---------------------------------------------------------------------------
%% intro.tex
%%
%% Introduction
%%
%% $Id: intro.tex 1477 2019-11-28 15:10:43Z lumurillo $
%% ---------------------------------------------------------------------------

\chapter{Introduction}
\label{chp:intro}

The diseases and pests control are becoming more relevant in the latest years due to the consequences that climate change brings, sucha as \cite{Coakley1999} the alteration of stages and rates of development of pathogens, the modification of host resistance and changes in the physiology of host-pathogen interactions. These issues cause changes in the geographical distribution and growth of plant species.

All these situations have the potential to cause catastrophic plant diseases, leading to the loss of food crops, which according to \cite{Strange2005} aggravates the deficit of food supply. In this sense, the tomato has a key role, due to its impact in the diet of people. As stated in \cite{Bergougnoux2014}, the tomato represents the most economically important vegetable crop worldwide due to the precense in the diet of millions of people, hence being one of the most produced vegetables after potatoes and before onions.

In \cite{Park2017} is mentioned the use of non-invasive methods to deal with crop disease. Normally, to combat this problem, the farmers use chemical agents as pesticides which has environmental consequences and health issues to the people that consume the vegetable.

In Costa Rica, the disease that causes the most damage to the tomato crops is the \emph{Phytophthora infestans} and the most important pest in the whiteflies \cite{MAG2007}.

In the development of automated applications for disease detection requires the use of a labeled dataset, which implies invest time consulting with experts for the labeling process. For that reason, this work explores methods to reduce this time as much as possible. The approach followed is the detection of anomalies where a set of data, in its majority, represents the healthy plants and the outliers can be seen as the affected plants or the anomalies.

The general goal of this project is the evaluation of different algorithms for anomaly detection, in the context of tomato images. The specific goals are the following:

\begin{itemize}
 \item To design a dataset to train the machine learning models.
 \item Selection of an anomaly detection algorithms.
 \item To evaluate the selected methods for detection anomaly in tomato images.
\end{itemize}

The rest of the work is structured as follows:

The state of the art chapter \ref{ch:marco} covers related works for the detection of anomalies in different contexts like in agriculture or health care; as well as additional concepts that underline this project.

In chapter \ref{ch:solution} the solution strategy to evaluate different deep learning architectures able to detect anomalies is described.

Chapter \ref{ch:results} presents a discussion of the performance of the different architecture and which of them are the best option for the detection of diseases in tomato.

Finally, the conclusions and the future work are summarized in chapter \ref{ch:conclusions}.

%%% Local Variables: 
%%% mode: latex
%%% TeX-master: "main"
%%% End: 

  \chapter{Literature review}
\label{ch:marco}

  \chapter{Proposed solution}
\label{ch:solution}

\section{Data preprocessing}

The data acquisition for this project has used a camera that captures in a resolution of 1920x1080. The tomato data was captured in tomato crops located in the Alajuela province of Costa Rica.

For the purpose of the architectures experiment was used 191520 images of a dimension fo 50x50 each. For this set of images, 80\% was used for testing and 20\% used for validation. This training validation data represents only healthy tomato plants. Also was used 798 images for testing and in this case, there are also unhealthy tomato photos.

Before using these images for the training of the models, a normalization was applied to convert the image values to be in a range between -1 to 1. Finally, some image data augmentation like rotation, zoom, width shift, among other image transformations was implemented.

\section{Architecture experiments}

In order to evaluate the different architectures for the anomaly detection in tomato plants, three experiments were proposed: one to explore the AnoGAN architecture, another to evaluate the spatial variational autoencoder and a third experiment to make use of the gaussian-mixture variational autoencoder.

All the experiment used the same dataset and for each architecture is tested its reconstruction of an input image. The expected behavior here is that, as the model is only trained with healthy images of tomato, if some input image presents a possible disease, the model shouldn't be able to reconstruct the affected area that reconstruction error or dissimilarity shall be an indication of the presence of an anomaly.

\section{Contribution}

The architecture proposed in this project is an adversarial anomaly detector inspired in the AnoGAN. This architecture is composed of the typical generator and discriminator of a typical GAN, but with the addition of an encoder \begin{math}E(x) = z\end{math} that is connected with the generator input.

The training process of this approach has two stages: First is the common training of a GAN model, but with just healthy images of tomato, and with the generation receiving as input random noise as latent variables. The second stage is the training of an encoder that shall be able to map to the latent space of the generator the necessary values to generate the most similar image from a query one.

In the training of the encoder, the generator and discriminator models are not trainables. Also, the hidden layers of the discriminator are used to extract features of the synthetic images and the input images. The goal is to train the encoder in a way that it makes that the generator creates images with similar features as the ones of the query image.

The proposed loss function for this anomaly detector model is the following:

\begin{equation}
 \mathcal{L}\left(\mathbf{E}(x)\right)=(1-\lambda) \cdot \mathcal{L}_{R}\left(\mathbf{E}(x)\right)+\lambda \cdot \mathcal{L}_{D}\left(\mathbf{E}(x)\right)
\end{equation}

where \begin{math}\lambda\end{math} represents the weights of the different termns of the function and \begin{math}\mathcal{L}_{R}\end{math} and \begin{math}\mathcal{L}_{D}\end{math} are defined as:

\begin{equation}
 \mathcal{L}_{R}\left(\mathbf{E}(x)\right)=\left\|\mathbf{x}-G\left(\mathbf{E}(x)\right)\right\|
\end{equation}

\begin{equation}
 \mathcal{L}_{D}\left(\mathbf{E}(x)\right)=\left\|\mathbf{f}(\mathbf{x})-\mathbf{f}\left(G\left(\mathbf{E}(x)\right)\right)\right\|
\end{equation}

where \begin{math}\mathbf{f}(x)\end{math} represents the feature extractor of the discriminator.

This approach overcomes one of the disadvantages of the AnoGAN that is its bad time performance. This modification of the AnoGAN architecture is similar to the one presented in \cite{Schlegl2019}.

  \chapter{Results and analysis}
\label{ch:results}

From the experiments of the three architectures: Adversarial Anomaly Detector, sVAE, and GM-VAE; the one that presented the best reconstruction in the Adversarial Anomaly Detector (see figures \ref{fig:anogan_rec}, \ref{fig:svae_rec} and \ref{fig:gmvae_rec}). For the cases of both autoencoders, the reconstructions suffer from lots of blurring making difficult to evaluate the presence of an anomaly and hence confirming one of the main problems this kind of autoencoders. The problems in the autoencoders could the provoked for the series of gaps the learned data distribution has due to the initial assumption of make use of simpler data distributions as a base.

\begin{figure}[H]
\begin{minipage}{\linewidth}
  \centering
  \begin{tabular}{ccc}
  \includegraphics[width=.26\linewidth]{anogan_sample1_original}
    & \includegraphics[width=.26\linewidth]{anogan_sample1_reconstruction} \\
  (a) & (b) \\
  \includegraphics[width=.26\linewidth]{anogan_sample2_original}
    & \includegraphics[width=.26\linewidth]{anogan_sample2_reconstruction} \\
  (c) & (d)
  \end{tabular}
  \end{minipage}
\caption[Reconstruction evaluation of the adversarial anomaly detector]{Reconstruction evaluation of the adversarial anomaly detector: a) origianl healty sample, b) reconstructed healty sample, c) original test sample and d) reconstructed test sample.}
\label{fig:anogan_rec}
\end{figure}

\begin{figure}[H]
\begin{minipage}{\linewidth}
  \centering
  \begin{tabular}{ccc}
  \includegraphics[width=.26\linewidth]{svae_sample1_original}
    & \includegraphics[width=.26\linewidth]{svae_sample1_reconstruction} \\
  (a) & (b) \\
  \includegraphics[width=.26\linewidth]{svae_sample2_original}
    & \includegraphics[width=.26\linewidth]{svae_sample2_reconstruction} \\
  (c) & (d)
  \end{tabular}
  \end{minipage}
\caption[Reconstruction evaluation of the sVAE]{Reconstruction evaluation of the sVAE: a) origianl healty sample, b) reconstructed healty sample, c) original test sample and d) reconstructed test sample.}
\label{fig:svae_rec}
\end{figure}

\begin{figure}[H]
\begin{minipage}{\linewidth}
  \centering
  \begin{tabular}{ccc}
  \includegraphics[width=.26\linewidth]{gmvae_sample1_original}
    & \includegraphics[width=.26\linewidth]{gmvae_sample1_reconstruction} \\
  (a) & (b) \\
  \includegraphics[width=.26\linewidth]{gmvae_sample2_original}
    & \includegraphics[width=.26\linewidth]{gmvae_sample2_reconstruction} \\
  (c) & (d)
  \end{tabular}
  \end{minipage}
\caption[Reconstruction evaluation of the GM-VAE]{Reconstruction evaluation of the GM-VAE: a) origianl healty sample, b) reconstructed healty sample, c) original test sample and d) reconstructed test sample.}
\label{fig:gmvae_rec}
\end{figure}

In the case of the AnoGAN, one of the more problematics parts was the training of the model, but with the available tomato date, it was possible to achieve the Nash equilibrium. The generated images are much better than the ones in the autoencoders. Its main disadvantage is the time needed to learn a set of latent variables to reconstruct some images. This problem is resolve with the introduction of an encoder in the AnoGAN architecture as shown in the table \ref{table:rec_time}.

\begin{table}[htb]
    \caption[Reconstruction time evaluation]{Reconstruction time evaluation.}
    \label{table:rec_time}
    \centering
    \begin{tabular}{ c c c }
        \hline
        Model & Reconstruction time (ms) \\
        \hline
        sVAE & 160.8 \\
        GM-VAE & 1383.1 \\
        AnoGAN & 6320.6 \\
        Adversarial Anomaly Detector & 255.3 \\
        \hline
    \end{tabular}
\end{table}

In figures \ref{fig:anogan_eval_test_image_1} and \ref{fig:anogan_eval_test_image_2} is possible to see the data distribution between normal and anomaly samples. For the case of the figure \ref{fig:anogan_eval_test_image_1} we can see a defined group of data that represents the healthy samples and a smaller amount of points around this healthy group that represents the anomaly data. In the figure \ref{fig:anogan_eval_test_image_2} we have the opposite case, where the amount of anomalies is bigger than the number of healthy samples.

\begin{figure}[htb]
  \centering
  \includegraphics[width=150mm]{anogan_t_sne1}
  \caption[Adversarial Anomaly Detector t-SNE evaluation of test image 1]{Adversarial Anomaly Detector t-SNE evaluation of test image 1}
  \label{fig:anogan_eval_test_image_1}
\end{figure}

\begin{figure}[htb]
  \centering
  \includegraphics[width=150mm]{anogan_t_sne2}
  \caption[Adversarial Anomaly Detector t-SNE evaluation of test image 2]{Adversarial Anomaly Detector t-SNE evaluation of test image 2}
  \label{fig:anogan_eval_test_image_2}
\end{figure}

  \chapter{Conclusions}
\label{ch:conclusions}

The experiments performed so far have shown a tendency of the variational autoencoder architectures to blur the reconstructed images. On the other side the AnoGAN has more promising results with better reconstruction images.

A modification to the AnoGAN architecture is proposed, allowing to considerably improve the reconstruction time.

As future work, the objective must be to improve the generator of the GAN. One possible strategy is in the implementation of a disentangling representation of the generator latent space \cite{Higgins2016}. Another possible approach is the use of a regularized in the latent space of the generator, in a similar way as proposed in \cite{Karras2018}. In order to validate the possible improvement in the generator, metrics that evaluate the GAN performance needs to be explored, along with metrics for the evaluation of the anomalies.

Finally, a segmentation process should be implemented in order to mark possible anomalous regions.


  %----------------------------------------------------------------------------
  % literature in bibtex way:
  % \bibliographystyle{sty/plainurl} % for english documents
  % \bibliography{literatura}
  % literature in biblatex/biber way
  \printbibliography[title={Bibliography},heading=bibintoc]
  %----------------------------------------------------------------------------

  %----------------------------------------------------------------------------
  \appendix
  %----------------------------------------------------------------------------

  \chapter{Demostración del teorema de Nyquist}

El título anterior es solo un ejemplo ilustrativo.  Éste teorema no ameritaría
un apéndice pues es parte normal del currículum de Electrónica, pero apéndices
usualmente involucran aspectos de esta índole, que se salen de la línea de la
tesis, pero que es conveniente incluir por completitud.

Los anexos contienen toda información adicional que se considere pertinente
agregar, como manuales de usuario, demostraciones matemáticas que se salen de
la línea principal de la tesis, pero que pueden considerarse parte de los
resultados del trabajo.


  %----------------------------------------------------------------------------
  \backmatter
  %----------------------------------------------------------------------------

  \printindex                % insert index into document. Don't forget to call
                             % "makeindex filename" first.
\end{document}
